\documentclass[11pt]{article}
\usepackage[a5paper,margin=16mm]{geometry}
\usepackage{fontspec}
\setmainfont{TeX Gyre Pagella}
\setsansfont{TeX Gyre Heros}
\usepackage{amsmath}
\usepackage{tikz}
\usepackage[most]{tcolorbox}
\tcbset{enhanced,fonttitle=\sffamily\bfseries,colback=blue!4,colframe=blue!50!black}
\newtcbtheorem[number within=section]{result}{Result}{colback=green!5,colframe=green!40!black}{res}
\newtcolorbox{callout}[1][]{colback=yellow!12,colframe=orange!70!black,title=Note,#1}

\begin{document}
\section{Boxes}
\begin{tcolorbox}[title=Plain box,sidebyside,lower separated=false]
Left side of the box.
\tcblower
Right side after the lower part.
\end{tcolorbox}

\begin{result}{Pythagoras}{pyth}
For a right triangle, $a^{2}+b^{2}=c^{2}$.
\end{result}
\begin{result}{Binomial}{bin}
$(x+y)^{2}=x^{2}+2xy+y^{2}$.
\end{result}
\begin{callout}
Results~\ref{res:pyth} and~\ref{res:bin} are numbered inside section~1.
\end{callout}

\section{A breakable box across pages}
\begin{tcolorbox}[breakable,title=Long box,colframe=purple!60!black,colback=purple!3]
\foreach \i in {1,...,22}{Line number \i{} of the breakable colour box, with enough text to wrap sometimes.\par}
\end{tcolorbox}
\end{document}
